\chapter{AES versus SM4: Key Lessons}
\begin{center}\begin{tabular}{P{.35\linewidth}P{.28\linewidth}P{.28\linewidth}}\toprule Property&AES-128&SM4\\\midrule Block size&128 bits&128 bits\\Key size&128 bits&128 bits\\Rounds&10&32\\State&4$\times$4 byte matrix&Four 32-bit words\\Nonlinearity&AES S-box&SM4 S-box\\Linear diffusion&MixColumns and ShiftRows&Rotations and XOR in \(L\)\\Last-round target&\(K_{10}\), 128 bits&\(rk_{31}\), 32 bits\\Master-key recovery&Reverse AES expansion&Reverse SM4 expansion from four consecutive round keys\\\bottomrule\end{tabular}\end{center}
\section{The same attack on AES and on SM4?}
\label{sec:aes_vs_sm4_attacks}
The \emph{principle} of each attack carries over to SM4 --- inject a fault before
a non-linear S-box layer, read it out through the (invertible) linear layer, and
filter key candidates with an S-box differential --- but the \emph{structure}
differs enough that SM4 is not simply a copy of the AES case.

\paragraph{Piret: AES versus SM4.}
\begin{center}\begin{tabular}{P{.26\linewidth}P{.32\linewidth}P{.32\linewidth}}\toprule
 & AES-128 Piret & SM4 Piret \\\midrule
Fault & one byte before a MixColumns & one byte at an S-box input \\
Linear layer exploited & MixColumns (mixes 4 bytes of a column) & \(L\) (mixes the 4 bytes of a 32-bit word) \\
How it is solved & match the MixColumns differential over a 4-byte key group & invert \(L\), then solve the S-box differential byte by byte \\
Faults needed (median) & \(\approx 2\) (round-8 model) & \(\approx 32\) \\
Recovered per campaign & all of \(K_{10}\) (128 bits) & one subkey \(rk_i\) (32 bits) \\\bottomrule
\end{tabular}\end{center}

\paragraph{Giraud: AES versus SM4.}
\begin{center}\begin{tabular}{P{.26\linewidth}P{.32\linewidth}P{.32\linewidth}}\toprule
 & AES-128 Giraud & SM4 Giraud \\\midrule
Fault & one bit before final SubBytes & one bit at an S-box input \\
Last-round linear layer & \textbf{none} (AES last round has no MixColumns) & \(L\) is still present \\
Output effect & exactly one ciphertext byte changes & a full 32-bit word changes \\
How it is solved & filter directly by Hamming weight one & invert \(L\) first, then filter by Hamming weight one \\
Faults needed (median) & 32 (known) / \(\approx 34.5\) (unknown) & 32 (known) / \(\approx 34.6\) (unknown) \\\bottomrule
\end{tabular}\end{center}

\paragraph{The key point.}
On \emph{AES}, Piret and Giraud are genuinely different: different round, byte
fault versus bit fault, and two faulty ciphertexts versus thirty-two. On
\emph{SM4} they move much closer together, which is probably why they can look
like ``the same thing''. The reason is structural: SM4 applies its linear layer
\(L\) in \emph{every} round, including the last, so both SM4 attacks must first
invert \(L\) and then solve the same S-box differential. What is left to
distinguish them is only the fault they assume --- a known whole-byte XOR for the
Piret model, or a single bit (known or unknown) for Giraud. In other words, on
SM4 the two attacks share one skeleton and differ only in the fault model, while
on AES the absence of MixColumns in the last round gives Giraud a shortcut that
Piret does not have.

\section{On SM4, what actually differs between Piret and Giraud?}
\label{sec:sm4_piret_vs_giraud}
Because SM4 keeps its linear layer \(L\) in every round, both attacks run the
\emph{same five steps}. Given a correct ciphertext and a faulty one for the same
input:
\begin{enumerate}
\item Inject the fault at an \textbf{S-box input} of the last round.
\item Invert the linear layer: \(\Delta B=L^{-1}(\Delta X_{35})=\tau(A)\xor\tau(A')\).
\item Solve the S-box differential byte by byte: keep the inputs \(a\) with
\(S_{\mathrm{SM4}}(a)\xor S_{\mathrm{SM4}}(a\xor e)=d\).
\item Convert the surviving \(a\) into bytes of the round key \(rk_{31}\).
\item Repeat for rounds 30, 29 and 28, then invert the key schedule.
\end{enumerate}
The implementation makes this literal: \texttt{giraud\_sm4.py} reuses the target
and the \(L^{-1}\) routine of \texttt{piret\_sm4.py}. The \emph{only} thing that
changes is the fault assumed in step~1, which in turn sets the value of \(e\)
used in step~3:

\begin{center}\begin{tabular}{P{.26\linewidth}P{.3\linewidth}P{.3\linewidth}}\toprule
 & SM4 Piret & SM4 Giraud \\\midrule
Fault injected & one \textbf{whole byte}, value known & one \textbf{single bit} \\
Mask \(e\) in \(S(a)\xor S(a\xor e)=d\) & a known byte, \hex{01}--\hex{FF} & a one-bit mask (\hex{01},\hex{02},\ldots,\hex{80}), known or not \\
Everything else (steps 2--5) & identical & identical \\\bottomrule
\end{tabular}\end{center}

\noindent The measured faulty-ciphertext counts confirm it: SM4 Piret and SM4
Giraud-known both need a median of \textbf{32} faults in the paired software dataset.
This observed cost is not a universal physical lower bound.
Only Giraud-\emph{unknown} rises slightly (median \(35\)) because
not knowing which bit flipped leaves eight possible values of \(e\) instead of
one, a weaker filter in step~3. In this project's SM4 adaptations the recovery
framework is shared; the fault restrictions and solver knowledge determine the
filter. This does not identify the two historical attacks as the same attack.
On AES the structures differ, because AES has no
MixColumns in the last round and Giraud exploits exactly that.

\paragraph{Cost of the master key.}
A further difference is how much a successful campaign buys. Reversing the AES
key schedule needs the single 128-bit last-round key \(K_{10}\), so one campaign
plus a schedule inversion recovers the master key. The SM4 schedule instead needs
\emph{four consecutive} 32-bit round keys (\(rk_{31},rk_{30},rk_{29},rk_{28}\)),
so the attacker must mount four successive campaigns --- faulting rounds 31, 30,
29 and 28 in turn and peeling off each recovered subkey --- before the master key
can be reconstructed. This difference follows from the round-key width and
schedule requirements; the round count alone does not establish attack cost.

\section{Why the analyses contain four SM4 configurations}
\label{sec:four_sm4_configurations}
The four plotted names describe experimental configurations, not four independently published attacks. The SM4 ``Piret'' label denotes this project's controlled known-mask DFA. The SM4 ``Giraud'' label denotes its single-bit-inspired filter. Neither name claims an unchanged transfer of an AES attack to SM4. In particular, the known XOR magnitude required by the SM4 example must not be attributed to every classical Piret--Quisquater fault model.

\begin{center}\small
\begin{tabular}{P{.23\linewidth}P{.31\linewidth}P{.34\linewidth}}\toprule
Plot label & Ordered masks used by the target & Information given to the solver\\\midrule
Piret: byte masks & 01, 02, 03, 04, 05, 06, 07, 08 & Exact mask for every observation; masks may contain several set bits.\\
Piret: matched bits & 01, 02, 04, 08, 10, 20, 40, 80 & Exact mask; the target uses only single-bit faults to match Giraud's sequence.\\
Giraud: known bit & 01, 02, 04, 08, 10, 20, 40, 80 & The exact flipped bit; a stronger-knowledge control experiment.\\
Giraud: unknown bit & Same single-bit sequence & Only the one-bit property. The simulator's actual mask is hidden from the candidate filter.\\\bottomrule
\end{tabular}\end{center}

\paragraph{Why two Piret configurations?}
The default byte-mask sequence includes multi-bit errors such as $\hex{03}=00000011_2$. The matched-bit configuration changes only the selected masks. It lets us compare candidate filters under identical observations rather than changing both the fault sequence and attacker knowledge simultaneously. A byte fault means corruption confined to a byte; it does not require all eight bits to change.

\paragraph{Why two Giraud configurations?}
The unknown-bit configuration preserves uncertainty about which bit flipped. The known-bit configuration supplies that information deliberately, to measure how much the uncertainty costs. Physical control of glitch timing does not automatically justify known-bit mode.

\subsection{The exact mathematical distinction}
Let $d$ be the observed S-box output difference after inverse diffusion and let $q$ be the observable XOR of the other three state-word bytes. For known mask $e$, the input and key candidate sets are
\[
\mathcal A_e(d)=\{a:S(a)\xor S(a\xor e)=d\},\qquad
\mathcal K_e(d,q)=\{a\xor q:a\in\mathcal A_e(d)\}.
\]
For an unknown flipped bit, define $\mathcal E=\{01,02,04,08,10,20,40,80\}_{16}$. The filter uses
\[
\mathcal K_{\rm unknown}(d,q)=\bigcup_{e\in\mathcal E}\mathcal K_e(d,q).
\]
The union is computed separately for each observation; it does not assume that the same unknown bit flipped in every encryption. Across observations the candidate sets are intersected. Since the actual single-bit mask is one member of $\mathcal E$, its known-mask candidate set is contained in the unknown-bit set for that observation.

For the first worked observation, $d=\hex{4F}$ and $q=\hex{07}$. Knowing $e=\hex{01}$ leaves key candidates $\{\hex{90},\hex{91}\}$. Hiding the bit leaves
\[
\{\hex{19},\hex{1B},\hex{6A},\hex{89},\hex{90},\hex{91},\hex{A0},\hex{B0},\hex{C9},\hex{EA}\}.
\]
Thus the two filters do not extract the same information from that pair.

\paragraph{Why some plotted results coincide.}
Piret with matched bits and Giraud with known bit test the same equation with the same observations. Their full candidate histories were checked for equality in all 100 paired trials. Separate graph panels expose these equal results. Equality here is a control result, not evidence that the historical AES attacks are identical.

\subsection{What the 100-trial comparison actually establishes}
All four configurations recovered and verified the full key in 100 seeded software trials. The three known-mask configurations consumed 32 faulty outputs each. Unknown-bit Giraud consumed a mean of 35.06, a median of 35 and a range of 32--40. These are dataset results, not universal fault requirements. Offline timing differences between equivalent filters describe implementation overhead.

With a maximum of two observations per target byte, unknown-bit Giraud reached $rk_{31}$ in 40\% of trials, $rk_{30}$ in 15\%, $rk_{29}$ in 7\%, and $rk_{28}$ and the master key in 3\%. The known-mask configurations completed all stages in these trials. Recovering the last subkey alone therefore must not be reported as full master-key recovery. Once all four consecutive tail words are unique, schedule inversion needs no further faulty ciphertexts.

\section*{Summary}
For AES Giraud, inject a one-bit fault before final-round SubBytes, test inverse-S-box differences for Hamming weight one, and intersect key-byte candidates. The corresponding SM4 adaptation first inverts the linear map $L$ to isolate the S-box differential.

For AES Piret--Quisquater, compare correct and faulty ciphertexts, invert the last S-box under candidate \(K_{10}\) bytes, match round-9 MixColumns differential patterns, intersect candidates, and reverse the AES key schedule.

For the illustrated SM4 model, compare correct and faulty ciphertexts, invert the linear map \(L\), test S-box input candidates using the known fault differences, derive \(rk_{31}\), and obtain further subkeys with additional appropriate experiments. Once four consecutive round keys are known, the SM4 key schedule can be reversed exactly.
